\documentclass[11pt, a4paper]{article}
\input{preamble.tex}
\switchlang{en}
\begin{document}

\title{1994 -- Elecom Lunaris Mouse}
\date{}
\maketitle
\selectlanguage{english}

The Lunaris mouse was introduced to the market in 1994 by Elecom, a Japanese electronics manufacturer which was founded eight years earlier \cite{elecom} and which by that time had already introduced two computer mice to the market. The Lunaris mouse was released in several variants, differing in connection interface and the presence or absence of the mouse resolution adjustment. This one (fig. \ref{luna}), model number M-L98MD, features an ADB interface and is therefore designed for Apple computers. In addition to this variant, models with PS/2, USB, and NEC PC-98 interfaces are also quite common.

\begin{figure}[h]
    \centering
    \includegraphics[scale=0.63]{1994_elecom_lunaris_mouse/pic_60.jpg}
    \caption{Elecom Lunaris Mouse} \label{luna}
\end{figure}

The mouse is clearly a designer product. This is no coincidence: its appearance was developed by Frog Design, the company of industrial designer Hartmut Esslinger \cite{frog}, the creator of the ``Snow White'' design language used in Apple's line of computers. Elecom's close attention to design was seen by the company as an advantage over competitors who focused mainly on technical specifications and price \cite{elecom_timeline}. The Lunaris product line, which began with this mouse and was later expanded to include other computer accessories, became the embodiment of this company policy in the 1990s.

Elecom's design efforts did not go unnoticed, and in 1994 the Lunaris mouse won Japan's prestigious Good Design Award \cite{award} (following the company's previous mouse model, the Elecom Graio 200 Mouse, which had won the same award the previous year).

\begin{figure}[h]
    \centering
    \includegraphics[scale=0.8]{1994_elecom_lunaris_mouse/top_30.jpg}
    \includegraphics[scale=0.8]{1994_elecom_lunaris_mouse/bottom_30.jpg}
    \caption{Elecom Lunaris Mouse, top and bottom views} \label{topandbottomluna}
\end{figure}

The body of the Lunaris mouse is made of beige plastic with a matte texture (there are also light or dark gray models; they can also be glossy). The mouse has a symmetrical shape, equally comfortable for left- and right-handed use (fig. \ref{topandbottomluna}). The mouse's design is based on a combination of clear geometric shapes (circle, arc, rectangle). The rotary switch ring is located at the base of the mouse body; it adjusts the mouse's resolution. Inside the ring on the underside, you'll find a small rubberized ball, a locking ring that allows you to clean the mouse, and a semicircular label with the mouse's technical specifications.

The top of the mouse body features several decorative longitudinal lines, the Lunaris product line logo, and two large buttons, which are embedded in the body's contours and separated by a longitudinal ridge (it was believed in the first half of the 1990s, that users would have difficulty distinguishing between the left and right buttons without a clear tactile cue). The body has a rectangular base with convex front and back sides. The diameter of the rotary switch ring exceeds the width of the rectangle, and the top of the ring is covered with protrusions in the form of segments of a sphere immersed in the volume of the main part of the body (fig. \ref{sizeluna}).

\begin{figure}[h]
    \centering
    \includegraphics[scale=0.45]{1994_elecom_lunaris_mouse/size_30.jpg}
    \caption{Elecom Lunaris Mouse on a graduated pad with a grid step of 1~cm} \label{sizeluna}
\end{figure}

Turning the switch counterclockwise decreases the mouse resolution (speeds up cursor movement), while turning it clockwise has the opposite effect. The ring provides step-by-step movement and is graduated from 50 to 1600 counts per inch. Lunaris mice with non-adjustable resolution retain the switch ring as a decorative element, without graduations; on the underside of such variants, the fixed CPI value of the mouse's resolution can be seen, unlike this model, which has the word ``variable'' in the corresponding place (fig. \ref{topandbottomluna}).

The convex, arched profile of the top of the mouse body and the elongated shape make it quite comfortable to use as a palm rest (fig. \ref{handluna}).

The Lunaris Mouse's motion tracking unit is a fairly typical optomechanical design from the early 1990s, with metal rollers and a black plastic shell covering the optocouplers (fig. \ref{insideluna}).

\begin{figure}[h!]
    \centering
    \includegraphics[scale=0.45]{1994_elecom_lunaris_mouse/hand_30.jpg}
    \caption{Elecom Lunaris Mouse with a human hand model} \label{handluna}
\end{figure}

The inverted position of the printed circuit board, raised on standoffs above the bottom of the case, is unusual, although not unique. In this case, it was used because of the rotary switch integrated into the bottom of the case.

\begin{figure}[h!]
    \centering
    \includegraphics[scale=0.83]{1994_elecom_lunaris_mouse/inside_60.jpg}
    \caption{Elecom Lunaris Mouse disassembled} \label{insideluna}
\end{figure}

\begin{thebibliography}{9}
    \bibitem {elecom} Elecom -- Wikipeda \url{https://en.wikipedia.org/wiki/Elecom}
    \bibitem {frog} Frog Design -- Wikipeda \url{https://en.wikipedia.org/wiki/Frog_Design}
    \bibitem{elecom_timeline} Elecom Sustainability Report 2025 \url{https://www.elecom.co.jp/ir/sustainability/pdf/i-2025_eng.pdf}
    \bibitem {award} LUNARIS-MOUSE M-L98MD, M-LMA, M-LP2 -- g-mark.org \url{https://www.g-mark.org/en/gallery/winners/9cdee7d5-803d-11ed-862b-0242ac130002}
\end{thebibliography}

\end{document}
